% !PRINTED 581
\weberpage{581}
\begin{center}
\textbf{\S\,156.\qquad \gesperrt{Transition to the singular moduli.}}
\end{center}

We now assume that $\omega$ satisfies a quadratic
equation:
\weq{1}{a\omega^2 + b\omega + c = 0}
whose \gesperrt{fundamental discriminant}
\weq{2}{\Delta = b^2 - 4ac}
is negative.

If then $j(\omega)$ is the invariant, the quantity $\lambda$ [§ 155,
(2)] satisfies the equation of degree 6:
\weq{3}{\lambda^6 - 9\cdot16\lambda^4 - [j(\omega)-27\cdot2^8]\lambda^2 + 64[j(\omega)-27\cdot64] = 0,}
whose coefficients are integral algebraic numbers. Consequently
$\lambda$ is itself an \gesperrt{integral algebraic number} (Bd. II, § 154, 11.)\footnote{It is this circumstance that particularly recommends the introduction of $\lambda$ in place of $\varkappa$, because $\varkappa$ is in general not an integer, but only $4\varkappa$ is.}
and in § 135 we have seen that $\lambda$ is a class invariant of the
discriminant $4\Delta$ or $16\Delta$.

We therefore adjoin the number $\lambda$ to the class field $\mathfrak{K}(\Delta)$
and obtain a class field:
\weq{4}{\begin{array}{l}
\mathfrak{L} = \mathfrak{K}(4\Delta), \text{ if } \Delta \equiv 0\ (\mathrm{mod}\ 4), \text{ or } \equiv 1\ (\mathrm{mod}\ 8).\\
\mathfrak{L} = \mathfrak{K}(16\Delta), \qquad \Delta \equiv 5\ (\mathrm{mod}\ 8).
\end{array}}
If $h$ is the class number of $\Delta$, and therefore the relative degree of $\mathfrak{K}(\Delta)$
with respect to the quadratic field $\Omega = \mathfrak{R}(\sqrt{\Delta})$, and $h'$ the
relative degree of $\mathfrak{L}$ with respect to $\Omega$, then (§ 100, § 123):
\weq{5}{\begin{array}{ll}
h' = h, & \text{if } \Delta \equiv 1\ (\mathrm{mod}\ 8),\\
h' = 2h, & \text{,,}\quad \Delta \equiv 0\ (\mathrm{mod}\ 4),\\
h' = 6h, & \text{,,}\quad \Delta \equiv 5\ (\mathrm{mod}\ 8).
\end{array}}
In the first of these three cases $\mathfrak{L}$ is therefore identical with $\mathfrak{K}(\Delta)$.

Through the substitution of the singular modulus $\omega$ the
coefficients $a_1, a_2 \ldots, b_1, b_2, \ldots$, in § 155, (13), (18) pass into integral
algebraic numbers, and from this it follows, Bd. II, § 154, 11.,
that the roots of $A(x)$, namely,
\weq{6}{\sqrt{\varkappa}\,\mathrm{sn}\,\Omega_{h,h'}, \quad \text{\gesperrt{are integers}},}
and that the roots of $B(x)$
\weq{7}{\sqrt{\varkappa}\,\frac{\mathrm{cn}\,\Omega_{h,h'}}{\mathrm{dn}\,\Omega_{h,h'}}, \quad \text{\gesperrt{are units}}.}
